\section{Conclusion}
\label{sec:conclusion}

CARAT validates candidate updates through hidden, class-balanced task certificates rather than relying only on geometric centrality or a single trusted reference direction. In the five-seed CIFAR-10/VGG19 and CIFAR-100/ResNet18 evaluation with $10\%$ malicious clients, CARAT attains an average rank of $1.75$--$2.75$ under the four conventional attacks and ranks first or second among defenses with Mimic results across all six dataset--heterogeneity settings. The CIFAR-100 results remain stable across seeds, while the CIFAR-10 results expose isolated chance-level failures under MinSum and MinMax as heterogeneity increases. Together with the strong unprotected Mean results on the conventional attacks and the small gaps in the separate $20\%$ component study, this evidence supports CARAT as a conditionally competitive reference-assisted mechanism within the evaluated threat level. Disjoint reference data, matched auxiliary-data budgets, adaptive attacks, and systematic evaluation at larger malicious-client fractions are important next steps.
